\clearpage
\section{Frobenius, subfield codes and higher multiplication gates}
\label{sec:frobenius-hierarchy-algebra}

A finite extension field supplies both a quantum register and arithmetic
operations on it.  The trace pairing identifies its position and momentum
labels over the prime field.  Relative trace then gives a logical momentum
coordinate for subfield codes and a Fourier-dual transfer between registers.
Multiplication in separate registers supplies reversible gates at every
higher Clifford level.  These statements hold in characteristic two as well
as at odd primes, with the choices of phase and computational basis explicit.

\subsection{Trace geometry and arithmetic Frobenius}

\begin{definition}[Trace-framed arithmetic register]
Fix a prime $p$ and a named primitive complex $p$-th root $\zeta_p$.
For a finite field $E/\mathbb F_p$ of cardinality $p^r$, the trace-framed
arithmetic register retains the field $E$, its Hilbert space
$\mathcal H_E=\ell^2(E)$ with orthonormal basis $|x\rangle$ indexed by
$x\in E$, and the trace character
\[
 \operatorname{Tr}_{E/\mathbb F_p}(x)=\sum_{j=0}^{r-1}x^{p^j},
 \qquad
 \psi_E(x)=\zeta_p^{\operatorname{Tr}_{E/\mathbb F_p}(x)}.
\]
Its prime-field phase space is $S_E=E\oplus E$ with form
\[
 \Omega_E((a,b),(a',b'))=
 \operatorname{Tr}_{E/\mathbb F_p}(ab'-a'b).
\]
The reference operators are
\[
 X_E(a)|x\rangle=|x+a\rangle,\qquad
 Z_E(b)|x\rangle=\psi_E(bx)|x\rangle,\qquad
 W_E(a,b)=Z_E(-b)X_E(a).
\]
An ordered list of registers has the tensor-product Hilbert space and
direct-sum prime-field phase space.  The empty list has Hilbert space
$\mathbb C$.  The field structures and ordered factors are retained data.
\end{definition}
\begin{scope}
The root of unity and reference operator convention are named choices.
No restriction beyond the stated finite-field hypotheses is imposed.
\end{scope}
\provenance{D3, D8, D1301}{theory/sidequests/frobenius-hierarchy/foundation.md}{theory/checks/frobenius\_hierarchy\_check.py (A1--A2)}{theory/verdicts/frobenius-hierarchy-algebra-r1.md}

\begin{definition}[Arithmetic Frobenius operator]
For a trace-framed arithmetic register with named $p,\zeta_p$, put
\[
 \sigma_E(x)=x^p,\qquad U_E|x\rangle=|\sigma_E(x)\rangle.
\]
For a list, Frobenius means the tensor product of these operators.
\end{definition}
\begin{scope}
This is the Frobenius of the named field structure.  It has the same field
as source and target.
\end{scope}
\provenance{D1302}{theory/sidequests/frobenius-hierarchy/foundation.md}{theory/checks/frobenius\_hierarchy\_check.py (A1)}{theory/verdicts/frobenius-hierarchy-algebra-r1.md}

\begin{definition}[Named embedding and transported relative trace]
For trace-framed arithmetic registers over the same named $p,\zeta_p$,
a named field embedding $i:K\to E$, $|K|=p^s$, and $r=ds$, define
\[
 T_i(x)=i^{-1}\left(\sum_{j=0}^{d-1}x^{p^{sj}}\right),\qquad
 \kappa_i=|\ker T_i|.
\]
Here $i^{-1}$ is restricted to $i(K)$.
For $j:L\to K$ the composite trace is denoted $T_{ij}:E\to L$.
\end{definition}
\begin{scope}
The trace formula taking values in $i(K)$ and its surjectivity are results
below, including when $p$ divides the extension degree.
\end{scope}
\provenance{D1303}{theory/sidequests/frobenius-hierarchy/foundation.md}{theory/checks/frobenius\_hierarchy\_check.py (A3)}{theory/verdicts/frobenius-hierarchy-algebra-r1.md}

\begin{theorem}[Relative trace and prime-field symplectic geometry]
{\normalfont\statusProved}\quad
For every named embedding $i:K\to E$ of finite fields over
$\mathbb F_p$, the transported relative trace is well defined,
surjective and $K$-linear, and
\[
 \kappa_i=\frac{|E|}{|K|},\qquad
 \operatorname{Tr}_{E/\mathbb F_p}(i(a)x)
 =\operatorname{Tr}_{K/\mathbb F_p}(aT_i(x)).
\]
For every tower $L\xrightarrow{j}K\xrightarrow{i}E$,
$T_{ij}=T_jT_i$.  The absolute trace pairing on $E$ is nondegenerate,
and $\Omega_E$ is alternating and nondegenerate over $\mathbb F_p$.
\end{theorem}
\begin{scope}
There is no division by the extension degree.  All prime characteristics
are included; no restriction beyond the stated hypotheses is imposed.
\end{scope}
\provenance{FRB-TRACE}{theory/sidequests/frobenius-hierarchy/foundation.md, section 1}{theory/checks/frobenius\_hierarchy\_check.py (A1--A3)}{theory/verdicts/frobenius-hierarchy-algebra-r1.md}
\begin{proof}
Put $q=|K|$ and $S(x)=\sum_{j=0}^{d-1}x^{q^j}$.
Since $x^{q^d}=x$, one has $S(x)^q=S(x)$; the $q$ roots of $t^q-t$
in $E$ are precisely $i(K)$.  The polynomial $S$ has degree
$q^{d-1}<q^d$ and nonzero leading coefficient, so it cannot vanish
on every point of $E$.  It is $i(K)$-linear and therefore surjects onto
$i(K)$.  Equal-size additive fibres give the kernel cardinality.
In a tower, expanding the two trace sums lists every Frobenius exponent
in the composite trace exactly once.  This proves transitivity and the
pairing identity.  If $c\ne0$, choose $u$ with absolute trace one and
put $x=c^{-1}u$; then $\operatorname{Tr}_{E/\mathbb F_p}(cx)=1$.
Thus the trace pairing is nondegenerate.  Pairing a vector $(a,b)$ with
$(0,x)$ and then $(x,0)$ proves nondegeneracy of $\Omega_E$;
$\Omega_E(v,v)=0$ directly in every characteristic.
\end{proof}

\begin{theorem}[Exact Frobenius covariance and atomic factorization]
{\normalfont\statusProved}\quad
For every trace-framed register $E/\mathbb F_p$, Frobenius is a field
automorphism, $U_E$ is unitary, both have order
$r=[E:\mathbb F_p]$, and
\[
 \Omega_E(\sigma_E v,\sigma_E w)=\Omega_E(v,w),\qquad
 U_EW_E(a,b)U_E^*=W_E(a^p,b^p).
\]
For a named $\mathbb F_p$-basis $(e_j)_{j=1}^r$ of $E$, there is a
unique trace-dual basis $(e^j)_{j=1}^r$.  If
$a=\sum_j a_je_j$ and $b=\sum_j b_je^j$, the coordinate unitary
$|\sum_jx_je_j\rangle\mapsto\bigotimes_j|x_j\rangle$ sends
$W_E(a,b)$ to $\bigotimes_jW_{\mathbb F_p}(a_j,b_j)$.
\end{theorem}
\begin{scope}
Position and momentum use dual bases.  No self-dual basis is required,
and no theorem about a lift of every symplectic map is asserted.
\end{scope}
\provenance{FRB-FROB}{theory/sidequests/frobenius-hierarchy/foundation.md, section 2}{theory/checks/frobenius\_hierarchy\_check.py (A1)}{theory/verdicts/frobenius-hierarchy-algebra-r1.md}
\begin{proof}
The absolute trace is invariant under the cyclic shift of Frobenius powers.
Consequently $U_EX_E(a)U_E^*=X_E(a^p)$ and
$U_EZ_E(b)U_E^*=Z_E(b^p)$, as is checked on the computational basis.
Multiplying in the reference order gives the exact Weyl formula.
The same trace invariance proves preservation of $\Omega_E$.
Frobenius has order $r$: its $r$-th power is the identity, while an
identity $x^{p^k}=x$ for all $x\in E$ with $0<k<r$ would contradict
the root bound for the polynomial $t^{p^k}-t$.
Nondegeneracy of the trace matrix supplies its inverse and hence the
dual basis.  Then
$\operatorname{Tr}_{E/\mathbb F_p}(b(x+a))=\sum_jb_j(x_j+a_j)$,
which factors the reference Weyl operator into the displayed tensor.
\end{proof}

\subsection{Subfield support and Fourier-dual transfers}

\begin{definition}[Inclusion and normalized trace-fibre transfers]
For a named embedding $i:K\to E$ of trace-framed registers, define
linear maps $J_i,V_i:\mathcal H_K\to\mathcal H_E$ by
\[
 J_i|a\rangle=|i(a)\rangle,\qquad
 V_i|a\rangle=\kappa_i^{-1/2}\sum_{T_i(x)=a}|x\rangle,
\]
using the positive real square root.  Their reverse arrows are their
Hilbert adjoints $J_i^*,V_i^*$.
\end{definition}
\begin{scope}
The source, target and positive normalization are fixed data; independent
phases on the transfers are not introduced.
\end{scope}
\provenance{D1304}{theory/sidequests/frobenius-hierarchy/transfers-example.md}{theory/checks/frobenius\_hierarchy\_check.py (A3)}{theory/verdicts/frobenius-hierarchy-algebra-r1.md}

\begin{definition}[Subfield support code and logical phase quotient]
For a named embedding $i:K\to E$ of trace-framed registers, put
\[
 \mathcal C_i=\operatorname{span}_{\mathbb C}
       \{|i(a)\rangle:a\in K\},\qquad P_i=J_iJ_i^*,\qquad
 N_i=\{0\}\oplus\ker T_i,\qquad
 G_i=\{Z_E(b):b\in\ker T_i\}.
\]
The proposed logical phase space is $N_i^\perp/N_i$, where
perpendicularity uses $\Omega_E$.  The logical coordinate map is
\[
 [(i(a),b)]\longmapsto (a,T_i(b))
 \quad\hbox{whenever }N_i^\perp=i(K)\oplus E.
\]
Its well-definedness and symplectic property are claims, not stipulations.
\end{definition}
\begin{scope}
The code is defined by computational-basis support in the embedded field.
It is not defined as the invariant-vector space of Frobenius.
\end{scope}
\provenance{D1305}{theory/sidequests/frobenius-hierarchy/foundation.md}{theory/checks/frobenius\_hierarchy\_check.py (A2)}{theory/verdicts/frobenius-hierarchy-algebra-r1.md}

\begin{theorem}[Subfield stabilizers and logical Weyl system]
{\normalfont\statusProved}\quad
For every named embedding $i:K\to E$ of trace-framed finite fields,
$\mathcal C_i$ is exactly the common $+1$ eigenspace of $G_i$, and
\[
 \dim\mathcal C_i=|K|,\qquad
 P_i=\frac1{\kappa_i}\sum_{b\in\ker T_i}Z_E(b),\qquad
 N_i^\perp=i(K)\oplus E.
\]
The logical coordinate map is a symplectic isomorphism
$N_i^\perp/N_i\cong S_K$.  Exactly the Weyl operators whose position
labels belong to $i(K)$ preserve the code, and
\[
 J_i^*W_E(i(a),b)J_i=W_K(a,T_i(b)).
\]
\end{theorem}
\begin{scope}
Logical momentum is a trace quotient.  The statement does not identify
it with inclusion of momentum labels from $K$.
\end{scope}
\provenance{FRB-CODE}{theory/sidequests/frobenius-hierarchy/foundation.md, section 3}{theory/checks/frobenius\_hierarchy\_check.py (A2)}{theory/verdicts/frobenius-hierarchy-algebra-r1.md}
\begin{proof}
The trace-pairing identity and nondegeneracy over $K$ show that the
annihilator of $i(K)$ is $\ker T_i$.  Taking annihilators again gives
$(\ker T_i)^\perp=i(K)$.  A basis vector is fixed by every phase
stabilizer precisely when its label belongs to this latter annihilator.
Character summation gives the displayed projector.  The same annihilator
calculation yields $N_i^\perp=i(K)\oplus E$; the displayed coordinate
map is surjective with kernel $N_i$ and preserves the form by trace
transitivity.  Finally, the Weyl operator sends a support coset to its
translate, and its phase on $|i(x)\rangle$ is
$\psi_K(-T_i(b)(x+a))$.  This proves the restriction formula.
\end{proof}

For an extension of degree $d$, relative trace gives $T_i(i(a))=da$.
Consequently $\psi_E(i(a))=\psi_K(a)^d$.  In particular, the included
momentum labels can all act trivially on the code when $p\mid d$.
The trace quotient remains valid because the full relative trace is
surjective in this case as well.

\begin{definition}[Negative-kernel arithmetic Fourier transform]
For a trace-framed register define
\[
 F_E|x\rangle=|E|^{-1/2}\sum_{y\in E}\psi_E(-xy)|y\rangle,
\]
with positive real square root.  A subscript on a tensor factor means
that this operator acts there and the identity acts on every other factor.
\end{definition}
\begin{scope}
The negative sign and positive normalization fix the Fourier convention.
No restriction beyond the finite-field hypotheses is imposed.
\end{scope}
\provenance{D1306}{theory/sidequests/frobenius-hierarchy/transfers-example.md}{theory/checks/frobenius\_hierarchy\_check.py (A3)}{theory/verdicts/frobenius-hierarchy-algebra-r1.md}

\begin{theorem}[Coherent inclusion and Fourier-dual trace transfer]
{\normalfont\statusProved}\quad
For every named embedding $i:K\to E$ of trace-framed finite fields,
$J_i,V_i$ are isometries, $F_E$ is unitary, and
\[
 F_EJ_i=V_iF_K,\qquad U_EJ_i=J_iU_K,\qquad U_EV_i=V_iU_K.
\]
For every tower $L\xrightarrow{j}K\xrightarrow{i}E$,
\[
 J_iJ_j=J_{ij},\qquad V_iV_j=V_{ij}.
\]
Identity embeddings give identity transfers.  Moreover,
\[
 F_EX_E(a)F_E^*=Z_E(-a),\qquad F_EZ_E(b)F_E^*=X_E(b).
\]
\end{theorem}
\begin{scope}
These are exact identities of maps with the displayed domains and
positive fibre normalizations.  Their adjoints reverse the domains.
\end{scope}
\provenance{FRB-TRANSFER}{theory/sidequests/frobenius-hierarchy/transfers-example.md, section 1}{theory/checks/frobenius\_hierarchy\_check.py (A3)}{theory/verdicts/frobenius-hierarchy-algebra-r1.md}
\begin{proof}
Different trace fibres are disjoint and each has $\kappa_i$ elements,
which proves the isometry assertion.  Character orthogonality proves
Fourier unitarity.  The coefficient of $|x\rangle$ on the two sides of
$F_EJ_i=V_iF_K$ is, respectively,
\[
 |E|^{-1/2}\psi_E(-i(a)x),\qquad
 (|K|\kappa_i)^{-1/2}\psi_K(-aT_i(x)).
\]
They agree by the trace-pairing identity and fibre cardinality.
In a tower, each output point occurs once, through its unique intermediate
trace label, and fibre cardinalities multiply.  Finally,
$T_i(x^p)=T_i(x)^p$, so Frobenius permutes corresponding fibres with
their normalization unchanged.  The two Fourier covariance formulas
follow by translating the kernel summation index.
\end{proof}

\subsection{Multiplication at exact higher Clifford levels}

\begin{definition}[Prime-field Pauli group and Clifford hierarchy]
For an ordered list $A=(E_1,\ldots,E_n)$ of trace-framed arithmetic
registers, let
\[
 \mathcal P_A=
 \left\{\lambda\bigotimes_jW_{E_j}(a_j,b_j):\lambda\in U(1)\right\}.
\]
Set $\mathcal C_1(A)=\mathcal P_A$ and, recursively for $k\geq1$,
\[
 \mathcal C_{k+1}(A)=
 \{U\text{ unitary}:UPU^*\in\mathcal C_k(A)
                         \text{ for every }P\in\mathcal P_A\}.
\]
The exact level of $U$ is the least positive $k$ with
$U\in\mathcal C_k(A)$, if such a $k$ exists.  All scalar phases are
included, in particular at $p=2$.  These sets are not stipulated to be
groups beyond the second level.
\end{definition}
\begin{scope}
The Pauli frame is the prime-field frame on all constituent atomic factors.
The finite group with phases restricted to $\mu_p$ is a different datum.
\end{scope}
\provenance{D1307}{theory/sidequests/frobenius-hierarchy/hierarchy.md}{theory/checks/frobenius\_hierarchy\_check.py (A5)}{theory/verdicts/frobenius-hierarchy-algebra-r1.md}

\begin{definition}[Reversible multiplication gate]
For a trace-framed register and integer $d\geq1$, define the linear
operator on $\mathcal H_E^{\otimes(d+1)}$ by
\[
 M_E^{(d)}|x_1,\ldots,x_d,z\rangle
 =\left|x_1,\ldots,x_d,z+\prod_{j=1}^dx_j\right\rangle.
\]
Each argument occupies a distinct register.
\end{definition}
\begin{scope}
Repeated powers of a variable within a single register are outside this
definition; the number of separate control registers is $d$.
\end{scope}
\provenance{D1308}{theory/sidequests/frobenius-hierarchy/hierarchy.md}{theory/checks/frobenius\_hierarchy\_check.py (A4--A5)}{theory/verdicts/frobenius-hierarchy-algebra-r1.md}

\begin{definition}[Phase polynomials and additive differences]
For an ordered list of trace-framed registers with configuration space
$A=\prod_jE_j$ and a function $f:A\to\mathbb F_p$, put
\[
 D_f|t\rangle=\zeta_p^{f(t)}|t\rangle,\qquad
 \delta_hf(t)=f(t)-f(t-h)\quad(h\in A).
\]
A polynomial of total degree at most $m$ means a polynomial expression
in the coordinates of named $\mathbb F_p$-bases, representing $f$, with
that degree bound.  The condition concerns existence of such an expression;
no unique polynomial representative is stipulated.
\end{definition}
\begin{scope}
The differences are operations on functions.  A polynomial expression can
have a larger degree than another expression defining the same function.
\end{scope}
\provenance{D1309}{theory/sidequests/frobenius-hierarchy/hierarchy.md}{theory/checks/frobenius\_hierarchy\_check.py (A5)}{theory/verdicts/frobenius-hierarchy-algebra-r1.md}

\begin{theorem}[Exact hierarchy level of reversible multiplication]
{\normalfont\statusProved}\quad
For every finite field $E/\mathbb F_p$ with named primitive phase and
every integer $d\geq1$, the reversible multiplication gate is unitary and
\[
 M_E^{(d)}\in\mathcal C_{d+1}(E,\ldots,E)
                   \setminus\mathcal C_d(E,\ldots,E).
\]
Thus its exact prime-field Clifford level is $d+1$.
\end{theorem}
\begin{scope}
Every characteristic and every positive $d$ is included.  The theorem
uses distinct registers and does not infer exact level from ordinary
polynomial degree alone.
\end{scope}
\provenance{FRB-HIERARCHY}{theory/sidequests/frobenius-hierarchy/hierarchy.md, sections 1--3}{theory/checks/frobenius\_hierarchy\_check.py (A5)}{theory/verdicts/frobenius-hierarchy-algebra-r1.md}
\begin{proof}
Subtracting the product gives the inverse permutation.  Fourier transform
on the target gives
\[
 F_{\mathrm{target}}M_E^{(d)}F_{\mathrm{target}}^*=D_f,
 \qquad f(x_1,\ldots,x_d,y)
       =-\operatorname{Tr}_{E/\mathbb F_p}\left(y\prod_jx_j\right).
\]
Fourier is Clifford, and Clifford conjugation preserves each hierarchy
level.  Multiplication on either side by a Pauli also preserves each
level; both facts follow inductively from the defining recursion.
The exact commutator identity is
\[
 D_fX(h)D_f^*X(h)^*=D_{\delta_hf}.
\]
A difference lowers the degree of a coordinate polynomial by at least
one.  Affine phases are Paulis by trace duality, so induction gives
$D_f\in\mathcal C_{d+1}$, since $f$ is multilinear in $d+1$ separate
arguments.  Conversely, a diagonal phase in $\mathcal C_k$ has every
$k$-fold additive difference constant: absorb the Pauli in the displayed
identity and induct down to a diagonal Pauli.

Now translate each of the $d$ distinct control registers by $1_E$.
The resulting $d$-fold difference is
$-\operatorname{Tr}_{E/\mathbb F_p}(y)$, which is nonconstant because
absolute trace is surjective.  Thus $D_f\notin\mathcal C_d$.
No factorial occurs in these separate-register differences, including
when $p\leq d$.  Fourier conjugation transfers both conclusions to
$M_E^{(d)}$.
\end{proof}

\begin{theorem}[Arithmetic covariance and encoded multiplication]
{\normalfont\statusProved}\quad
For every named embedding $i:K\to E$ of trace-framed finite fields
and every $d\geq1$, multiplication commutes with simultaneous Frobenius
and satisfies
\[
 U_E^{\otimes(d+1)}M_E^{(d)}
 =M_E^{(d)}U_E^{\otimes(d+1)},\qquad
 M_E^{(d)}J_i^{\otimes(d+1)}
 =J_i^{\otimes(d+1)}M_K^{(d)}.
\]
The latter identity also holds for inverse gates.  Multiplication
preserves $\mathcal C_i^{\otimes(d+1)}$; its encoded logical action
is $M_K^{(d)}$ and has exact level $d+1$ for the logical prime-field
Pauli frame supplied by the trace quotient.
\end{theorem}
\begin{scope}
The intertwining uses the support encoding $J_i$.  No corresponding
multiplicative intertwining is asserted for trace-fibre encoding $V_i$.
\end{scope}
\provenance{FRB-NATURAL}{theory/sidequests/frobenius-hierarchy/hierarchy.md, section 4}{theory/checks/frobenius\_hierarchy\_check.py (A4)}{theory/verdicts/frobenius-hierarchy-algebra-r1.md}
\begin{proof}
Frobenius and field embeddings preserve sums and products, so the two
orders of each composition agree on every computational basis point.
Compressing by the tensor encoding's adjoint gives the logical gate.
The trace quotient identifies its Pauli frame with that over $K$, to
which the exact-level theorem applies.
\end{proof}

\begin{proposition}[Coherent multiplication return probability]
{\normalfont\statusProved}\quad
For a trace-framed finite field $E$, put $Q=|E|$ and
\[
 |+_E\rangle=Q^{-1/2}\sum_x|x\rangle,\qquad
 D=F_{\mathrm{target}}M_E^{(2)}F_{\mathrm{target}}^*.
\]
The return amplitude of $D$ on $|+_E\rangle^{\otimes3}$ is
$(2Q-1)/Q^2$, and its return probability is $(2Q-1)^2/Q^4$.
\end{proposition}
\begin{scope}
The input, unitary and return effect are fixed as stated.  The rational
formula alone does not supply a specialization functor or reference trace.
\end{scope}
\provenance{FRB-NATURAL}{theory/sidequests/frobenius-hierarchy/hierarchy.md, section 5}{theory/checks/frobenius\_hierarchy\_check.py (A5-interference)}{theory/verdicts/frobenius-hierarchy-algebra-r1.md}
\begin{proof}
The amplitude is
\[
 \frac1{Q^3}\sum_{x,y,t\in E}\psi_E(-txy)
 =\frac{2Q-1}{Q^2},
\]
because summing over $t$ gives $Q$ exactly when $xy=0$ and zero
otherwise.  There are $2Q-1$ pairs with $xy=0$.  The return probability
is $(2Q-1)^2/Q^4$, giving $9/16$ at $Q=2$, whereas identity evolution
gives one.  The rational function takes the value one at $Q=1$; that
calculation does not supply a specialization functor or its reference trace.
\end{proof}

\subsection{The binary degree-four tower}

\begin{definition}[Binary degree-four field tower]
Put
\[
 K=\mathbb F_2[a]/(a^2+a+1),\qquad
 E=K[b]/(b^2+b+a).
\]
The named embeddings are the constant embeddings
$\mathbb F_2\to K\to E$.  Write an element of $E$ as
$(u,v)=u+vb$, with $u,v\in K$, and use $\zeta_2=-1$.
Irreducibility and the resulting field orders are claims.
\end{definition}
\begin{scope}
The polynomial models and embeddings fix concrete coordinates for the
example.  No basis-independent coordinate decomposition is asserted.
\end{scope}
\provenance{D1310}{theory/sidequests/frobenius-hierarchy/transfers-example.md}{theory/checks/frobenius\_hierarchy\_check.py (A6)}{theory/verdicts/frobenius-hierarchy-algebra-r1.md}

\begin{theorem}[Frobenius, code and transfer in the binary quartic field]
{\normalfont\statusProved}\quad
For the displayed binary tower, $K$ and $E$ are fields of cardinalities
four and sixteen.  For $u,v\in K$,
\[
 \sigma_E(u,v)=(u^2+av^2,v^2),\qquad
 \sigma_E^2(u,v)=(u+v,v),\qquad
 \sigma_E^4=\mathrm{id},\qquad \operatorname{ord}(\sigma_E)=4.
\]
For $i:K\to E$, the relative and absolute traces are
\[
 T_i(u,v)=v,\qquad
 \operatorname{Tr}_{E/\mathbb F_2}(u,v)=v+v^2.
\]
The subfield code has basis $|u,0\rangle$, dimension four, and transfers
\[
 J_i|u\rangle=|u,0\rangle,\qquad
 V_i|v\rangle=\frac12\sum_{u\in K}|u,v\rangle.
\]
The full trace transfer from $\mathbb F_2$ has coefficient $1/\sqrt8$
on each absolute-trace fibre.  The fixed-vector space of $U_E^2$ has
dimension ten, whereas the support code has dimension four.
\end{theorem}
\begin{scope}
The two different invariant spaces are computed in these named
coordinates; this does not identify Hilbert-space fixed vectors with
arithmetic fixed field elements in general.
\end{scope}
\provenance{FRB-EXAMPLE}{theory/sidequests/frobenius-hierarchy/transfers-example.md, section 3}{theory/checks/frobenius\_hierarchy\_check.py (A6)}{theory/verdicts/frobenius-hierarchy-algebra-r1.md}
\begin{proof}
The first quadratic has no root in $\mathbb F_2$.
In $K$, the values of $x^2+x$ are $0,0,1,1$ at
$0,1,a,a+1$, so the second quadratic has no root either.
Using $b^2=b+a$ gives the first Frobenius formula.  Squaring again
and using $a^2+a=1$ gives the second, which is nonidentity and squares
to the identity.  The relative trace $x+x^4$ is therefore $(v,0)$,
giving the trace and fibre formulas.  Logical position $k$ may be
represented by $(k,0)$ and logical momentum $m$ by $(0,m)$; the
momentums $(m,0)$ are phase stabilizers.  Finally $U_E^2$ has four
singleton basis orbits and six orbits of length two.  A fixed vector
has constant coefficients on each orbit, so its fixed space has
dimension ten.
\end{proof}

The tower also supports the two-control multiplication gate over $K$
inside that over $E$: both physical and logical gates have exact level
three, and the displayed Frobenius commutes with them on their respective
registers.  The field multiplication, Frobenius permutation, support code
and trace fibres remain explicitly available to the process construction.
